\documentclass[11pt]{article}

\usepackage[T1]{fontenc}
\usepackage[utf8]{inputenc}
\usepackage{lmodern}
\usepackage[a4paper,margin=3cm]{geometry}
\usepackage{microtype}
\usepackage{parskip}

\pagestyle{empty}

\begin{document}

\begin{center}
{\LARGE A Town Gate That Learns to Wait}\\[0.6em]
{\large A plain-language account of the ``stigmergic-quorum admission
control'' technique}\\[0.3em]
{\small \texttt{moderation-pipeline-prototype} --- August 2026}
\end{center}

\vspace{1em}

\textbf{The problem.} Some disputes that reach a moderation system are
genuinely hard to settle: two sides each have arguments that attack the
other's, and the easy, fast method for weighing them comes back with no
clear answer. When that happens, the system falls back to a much more
thorough, much more expensive method of sorting the whole tangle out ---
worth doing, but not something you want to do for free, on demand,
for anyone who asks. The trouble is in how that expensive method gets
triggered. In the original design, the very first person who shows up
with a plausible-looking argument gets the case waved through to the
expensive process, instantly. That is exactly the kind of rule an
attacker loves: show up once, in bulk, with a pile of manufactured
arguments, and the expensive gate opens for free, as many times as you
like.

\textbf{The idea.} The fix borrows two tricks from nature, and it is
worth taking them one at a time, because each solves a different half
of the problem.

The first trick is how ant colonies remember which path to food is
worth following. An ant does not vote once and settle the matter; it
lays down a scent as it walks, and that scent fades steadily unless
other ants, arriving independently, walk the same way and refresh it. A
path only stays marked if it keeps earning fresh footsteps from
different ants --- one ant pacing back and forth on its own trail all
day leaves a trail that is, for practical purposes, exactly as strong as
one visit. Applied here: every disputed case carries a score that fades
a little every tick of the clock, and a new argument only tops the score
back up by very little if it comes from a source that has already
weighed in recently. Genuinely new, independent concern keeps a
case's score alive; the same complainer shouting the same thing over and
over does not.

The second trick is how bacteria decide, as a colony, when there are
enough of them to act together --- a mechanism called quorum sensing.
No single bacterium makes that call alone; the whole colony reads a
shared chemical signal that reflects how crowded things already are, and
raises or lowers its collective threshold for action accordingly.
Applied here: the bar a case must clear to reach the expensive process
is not fixed. It rises automatically whenever that expensive process is
already busy with other cases, exactly the way a restaurant becomes
harder to get a table at as it fills up, without anyone at the door
consciously deciding to be stricter.

\textbf{Put together.} A disputed case's score decays like an ant trail
losing its scent, is topped up only by reinforcement that looks like it
came from someone new, is capped so no single moment of activity can
push it past a sensible cap, and must clear a bar that rises with how
busy the expensive process already is. A single burst of manufactured
outrage, however large, cannot move a decaying, capped score past a
threshold in one instant --- there simply is not enough weight allowed
into any one tick of the clock, no matter how many fake voices show up
at once.

\textbf{What was actually measured, not just hoped for.} The mechanism
was tested against its own predecessor --- the instant-admission rule
--- under several kinds of attack, and the two were compared honestly,
including on the case where the new mechanism does \emph{not} fully
win. A single flood, of any size from one voice to five thousand,
landing in one moment: fully blocked, every time, where the old rule let
every single one through. The same flood sustained for a very long time
from one traceable source: still fully blocked, provably, not just
``blocked so far in the tests that were run.'' Two hundred separate
disputes each hit with a one-shot burst, alongside fifteen genuine,
organically-growing disputes running at the same time: the old rule let
every attack through and admitted every genuine case too, at a real
processing cost of 215 cases; the new mechanism let through zero of the
attacks while still admitting every one of the genuine fifteen, at a
real processing cost of 15 --- ninety-three percent less work forced
onto the expensive process. The honest exception, stated plainly rather
than buried: a patient attacker willing to manufacture a fresh,
never-reused, independent-looking identity every single tick of the
clock, for several ticks in a row, is not fully stopped. That attacker
still gets through eventually, at roughly five times the effort and time
the old rule would have cost them. The mechanism turns a free attack
into an expensive one; it does not turn it into an impossible one, and
that limit is reported here rather than left for someone else to
discover the hard way.

\textbf{What it deliberately does not do.} It says nothing about who is
\emph{right} in the underlying dispute --- that judgment is left entirely
to the expensive process it is guarding, once a case is let through. It
assumes there is some real cost, even a small one, to appearing as a
brand-new independent voice every single tick; if that assumption ever
stops holding for a given attacker, the delay is the whole defense, not
a backup behind something stronger. And the fading-and-capping trick and
the rising-bar trick, taken completely apart from each other, are each
individually well known outside biology too --- the first resembles an
ordinary rate limiter used all over computing, the second resembles a
decades-old networking technique for easing congestion before a queue
overflows. What a dedicated search of the relevant research literature
did not turn up is this particular combination of all three pieces,
guarding specifically this kind of expensive, hard-to-parallelize
fallback process. That is the narrow claim being made, and only to the
extent that search was thorough --- not a claim that no one, anywhere,
has ever combined these ideas before.

\textbf{Why it is worth writing down.} Nothing about the expensive
process itself needed to change; the fix sits entirely at the door in
front of it. The individual ingredients --- a fading score, a cap on
how much one moment can contribute, a bar that rises with crowding ---
are each borrowed from places with decades of their own track record,
biological and computational alike. What is reported here is how they
compose, tested honestly against the attack it was built to resist, with
the case where it only partly succeeds stated as plainly as the cases
where it fully does.

\end{document}
